% problem-000063 5 一氧化氮氢化动力学
\section*{5. 一氧化氮氢化动力学}
$
2 \mathrm{NO} (\mathrm{g}) + 2 \mathrm{H} _ {2} (\mathrm{g}) \rightarrow \mathrm{N} _ {2} (\mathrm{g}) + 2 \mathrm{H} _ {2} \mathrm{O} (\mathrm{g})
$

实验得到了起始反应速率的⼀组数据如下：

\textit{（原卷此处含表格/仪器试剂表，详见 PDF 或 MD 源文件。）}

\subsection*{5-1}
根据以上实验结果，求出反应物的反应级数（取整数），并计算反应的平均速率常数。

\subsection*{5-2}
NO氢化还原反应的动⼒学研究发现：

(a) NO为含有奇数电⼦的分⼦，易于发⽣聚合，并达到平衡。

(b) 反应中间产物除了NO聚合产物外，还有其他氮氧化物。

(c) NO并⾮直接被还原变成产物 $\mathrm { N } _ { 2 }$ 及 $\mathrm { H _ { 2 } O } ,$ ，⽽是通过中间产物的还原实现的。

\subsection*{5-2}
-1 根据上述事实，提出NO氢化还原反应的机理。

\subsection*{5-2}
-2 根据提出的反应机理（紧接平衡反应的步骤为速控步骤），推导反应速率⽅程，并与实验测定的反应速率⽅程进⾏⽐较，说明反应机理的合理性。

\subsection*{5-2}
-3 根据推导的反应速率⽅程，推算反应表观活化能与相关基元步骤活化能之间的关系。

$5 – 2 – 4 \mathrm { H N _ { 3 } }$ 是⾮直线型分⼦，NO能够取代 $\mathrm { H N } _ { 3 }$ 中的H形成 $\mathrm { N _ { 4 } O _ { : } }$ ，画出 $\mathrm { N } _ { 4 } \mathrm { O }$ 分⼦最稳定的结构式，并按N-O键键⻓由⻓到短的顺序排列NO、NO+和 $\mathrm { N } _ { 4 } \mathrm { O } _ { \circ }$
